% =====================================================================
%  capitulos/muestra.tex — Capítulo de muestra del diseño
%  Fragmentos del cap. 2 usados solo para mostrar cada componente.
% =====================================================================
\capitulo{Las reglas de la gramática}{Augusto M. Trombetta}

\section{Introducción}

Las reglas de la gramática no parecen ser todas iguales: unas definen las
características constitutivas propias de los objetos de una lengua, otras
establecen normas de corrección y etiqueta propias del uso social de esa
lengua. En un caso, las reglas se descubren (por ejemplo, \forma{la oración
consta de un sujeto y un predicado}); en el otro, las reglas se fijan (por
ejemplo, \forma{no debe haber coma entre el sujeto y el predicado}).\footnote{En
  su ensayo de 1969, \textit{Actos de habla}, Searle retoma su distinción entre
  reglas regulativas y constitutivas \parencite[48--49]{searle1994}.}

\section{Concepciones de la gramática}

La valoración de una combinación de la lengua como \termino{correcta} o
\termino{incorrecta} solo adquiere sentido en el marco de una gramática
normativa. En contraste, una gramática descriptiva valora las estructuras como
\termino{gramaticales} o \termino{agramaticales} \parencite[29]{bosque2009}:

\begin{table}[htbp]
  \centering
  \begin{tabularx}{\linewidth}{@{}>{\itshape}l C C@{}}
    \toprule
    Concepciones de la gramática & Normativa o prescriptiva & Descriptiva o científica \\
    \midrule
    Tipos de reglas          & Regulativas (externas)  & Constitutivas (internas) \\
    Criterios de evaluación  & Corrección/Incorrección & Gramaticalidad/\hspace{0pt}Agramaticalidad \\
    \bottomrule
  \end{tabularx}
  \caption{Concepciones de la gramática}
  \label{tab:concepciones}
\end{table}

\begin{ejercicio}
Considere las siguientes definiciones, representativas de las dos concepciones
señaladas sobre la gramática, y complete con los valores «\mas» o «\menos» el
cuadro presentado más abajo, como se muestra en el ejemplo resuelto~(1).

\begin{exe}
  \ex\label{ej1:salva} La \textit{gramática de la actual lengua castellana} no es
      otra cosa que \textit{el conjunto ordenado de las reglas de lenguaje que
      vemos observadas en los escritos o conversación de las personas doctas}
      \parencite[1]{salva1852}.
  \ex La gramática de L será, pues, un ingenio que genere todas las secuencias
      gramaticales de L y ninguna de las agramaticales \parencite[27]{chomsky1974a}.
\end{exe}

\begin{tabularx}{\linewidth}{@{}L Y{1.2cm} Y{1.2cm}@{}}
  \toprule
  De acuerdo con las definiciones dadas, la gramática de una lengua\ldots & 1 & 2 \\
  \midrule
  \ldots\ remite al uso lingüístico de ciertos grupos sociales.            & \mas   & \\
  \ldots\ constituye un mecanismo capaz de producir estructuras.          & \menos & \\
  \ldots\ conforma una disciplina normativa o prescriptiva.               & \mas   & \\
  \bottomrule
\end{tabularx}

\textit{Nota}: la columna correspondiente a 1 muestra el ejemplo resuelto.
\end{ejercicio}

\section{Ambigüedad}

Una secuencia posee ambigüedad estructural «cuando puede corresponderle más de
una estructura sintáctica, y por tanto más de una interpretación»
\parencite[35]{bosque2009}. Una oración como \exref{ex:secarropa} admite dos
lecturas:

\begin{exe}
  \ex\label{ex:secarropa} Por favor, cuando termine de usar el secarropa, limpie el
      filtro de pelusas.
\end{exe}

En \exref{ex:secarropa} se puede considerar que [\forma{el filtro}] y
[\forma{de pelusas}] son dos modificadores independientes del verbo
\exref{ex:lectura1}, o bien que [\forma{de pelusas}] modifica a [\forma{el
filtro}] \exref{ex:lectura2}:

\begin{exe}
  \ex\label{ex:lectura1}
    \begin{xlist}
      \ex Por favor, cuando termine de usar el secarropa, limpie [el filtro] [de pelusas].
      \ex Por favor, cuando termine de usar el secarropa, limpie [de pelusas] [el filtro].
    \end{xlist}
  \ex\label{ex:lectura2}
    \begin{xlist}
      \ex Por favor, cuando termine de usar el secarropa, limpie [el filtro de pelusas].
      \ex[*]{Por favor, cuando termine de usar el secarropa, limpie [de pelusas el filtro].}
    \end{xlist}
\end{exe}

La estructura de \exref{ex:lectura2} puede representarse así:

\begin{arbol}
  [SV [V [limpie]] [SD [D [el]] [SN [N [filtro]] [SP [de pelusas, roof]]]]]
\end{arbol}

Ejemplo glosado (Leipzig), con duplicación de clítico rioplatense:

\begin{exe}
  \ex\label{ex:glosa}
    \gll Lo vi a Juan.\\
         \Acc.3\Sg.\M{} ver.\Pst.1\Sg{} \Dom{} Juan\\
    \glt ‘Vi a Juan.’
\end{exe}

\begin{ejercicio}[Ambigüedad]
Considere las siguientes oraciones y conteste las preguntas que se formulan
abajo, como se muestra en el ejemplo resuelto~(1).

\begin{exe}
  \ex\label{ej7:escultura} Me trajeron esta escultura de Francia.
  \ex Antes de salir a navegar, acomodaron correctamente las velas.
  \ex Yo no recibiría a tu hermano en mi oficina.
\end{exe}

\begin{consignas}
  \item Las tres oraciones pueden tener más de una interpretación. Realice
        paráfrasis que pongan de relieve los diferentes significados.
  \item ¿Observa alguna mala formación de estructura gramatical?
\end{consignas}

\begin{resuelto}
\ref{ej7:escultura} (a) La oración está bien formada y es gramatical, aunque
resulte ambigua: [Me trajeron [esta escultura [de Francia]]] / [Me trajeron
[esta escultura] [de Francia]]. (b) La ambigüedad se podría resolver
reestructurando el orden de los constituyentes.
\end{resuelto}
\end{ejercicio}

Después del ejercicio, la numeración del texto continúa:

\begin{exe}
  \ex Si nos matamos, nos matamos.
\end{exe}

\begin{recuadro}[Reflexión final]
Las tipologías referidas a las diferentes variedades lingüísticas coexisten en
cada individuo que articula una oración en su variedad de lengua, una variedad
que es a la vez temporal, geográfica y social, y que se produce en un
determinado contexto situacional.
\end{recuadro}

\begin{destacado}
\textbf{Gramaticalidad} y \textbf{aceptabilidad} son nociones distintas: la
primera pertenece a la teoría de la competencia; la segunda, a la de la
actuación.
\end{destacado}

% --- Secciones finales -------------------------------------------------
\lecturasrecomendadas

\solucionario

\begin{solucion}{1}
\begin{tabularx}{\linewidth}{@{}L Y{1.2cm}@{}}
  \toprule
  De acuerdo con las definiciones dadas, la gramática de una lengua\ldots & 2 \\
  \midrule
  \ldots\ constituye un mecanismo capaz de producir estructuras. & \textbf{\mas} \\
  \bottomrule
\end{tabularx}
\end{solucion}

\begin{solucion}{2}
Las tres oraciones del ejercicio están bien formadas, aunque son ambiguas. En
(2), la ambigüedad es léxica; en (3), sintáctica.
\end{solucion}

